%% The example numbers below are those cited in the corresponding module
%% docstrings from The asymmetric fundamental theorem of matrix product states.
%% That lecture note is not included in Papers/ or Notes/; the statements here
%% have been compared with the linked declarations, not with an unseen source.

\section{Examples of multi-block compressions}

The following examples distinguish irreducibility from normality, equality of
word traces from gauge equivalence, and agreement of reading maps on matched
subspaces from agreement on the whole bond space. The example numbers refer to
\emph{The asymmetric fundamental theorem of matrix product states}. Coordinates
are indexed from $0$, and $E_{ij}$ and $e_i$ denote the matrix units and
coordinate vectors.

\subsection{Irreducibility and word traces}

\begin{lemma}[An irreducible periodic tensor]
    % Lecture note: equation (1.1).
    \label{lem:asym_example_periodic}
    \lean{MultiBlockExamples.Periodic.tensorC,
        MultiBlockExamples.Periodic.isIrreducibleAction_tensorC,
        MultiBlockExamples.Periodic.wordSpan_tensorC_of_even,
        MultiBlockExamples.Periodic.wordSpan_tensorC_of_odd,
        MultiBlockExamples.Periodic.finrank_wordSpan_tensorC,
        MultiBlockExamples.Periodic.not_isNormal_tensorC}
    \leanok
    The tensor of equation~(1.1) has $C^0=E_{01}$ and $C^1=E_{10}$ in $\MN{2}$.
    The only common invariant
    subspaces are $0$ and $\C^2$. For every positive length $n$, the span of the
    words of length $n$ is
    \begin{align}
        \spn_{\C}\{C^w:|w|=n\}
         & =\begin{cases}
                \spn_{\C}\{E_{00},E_{11}\}, & n\text{ even}, \\
                \spn_{\C}\{E_{01},E_{10}\}, & n\text{ odd}.
            \end{cases}
    \end{align}
    Each span has dimension $2$, so $C$ is irreducible but not normal.
\end{lemma}

\begin{proof}
    \leanok
    Multiplication by either letter changes the parity of the nonzero matrix
    entries. The alternating words give the two indicated matrix units at
    every positive length. If a common invariant subspace contains a nonzero
    vector, applying the two letters produces both coordinate vectors.
\end{proof}

\begin{lemma}[Equal word traces without a gauge]
    % Lecture note: Example 4.5.
    \label{lem:asym_example_no_gauge}
    \lean{MultiBlockExamples.NoGauge.tensorB,
        MultiBlockExamples.NoGauge.tensorB',
        MultiBlockExamples.NoGauge.evalWord_tensorB',
        MultiBlockExamples.NoGauge.evalWord_tensorB_apply,
        MultiBlockExamples.NoGauge.trace_evalWord_tensorB_eq,
        MultiBlockExamples.NoGauge.trace_evalWord_tensorB'_eq_one_of_eq_replicate,
        MultiBlockExamples.NoGauge.trace_evalWord_tensorB'_eq_zero_of_ne_replicate,
        MultiBlockExamples.NoGauge.tensorB_one_mul_zero_and_tensorB'_commute,
        MultiBlockExamples.NoGauge.isInvariantSubmodule_span_tensorB_iff,
        MultiBlockExamples.NoGauge.isInvariantSubmodule_span_tensorB'_iff,
        MultiBlockExamples.NoGauge.not_exists_conj_tensorB_eq_tensorB'}
    \leanok
    In Example~4.5, put $B=(E_{00}+E_{01},E_{11})$ and
    $B'=(E_{00},E_{11})$.
    For every word $w$, $\tr(B^w)=\tr((B')^w)$. This common trace is $2$ for
    the empty word, $1$ for a nonempty word using only one letter, and $0$ for
    a word using both letters. Nevertheless, no invertible $X$ satisfies
    $XB^iX^{-1}=(B')^i$ for both letters. The only common invariant line of
    $B$ is $\C e_0$, whereas those of $B'$ are $\C e_0$ and $\C e_1$.
\end{lemma}

\begin{proof}
    \leanok
    The matrix $(B')^w$ is diagonal, with entry $1$ at $k$ precisely when every
    letter of $w$ equals $k$. The product $B^w$ is upper triangular with the
    same diagonal. Also $B^0B^1-B^1B^0=E_{01}$, while the letters of $B'$ commute;
    conjugation preserves commutation. Solving $B^iv=\lambda_i v$ and
    $(B')^iv=\lambda_i v$ gives the stated invariant lines.
\end{proof}

\begin{lemma}[Interpolation without a constant term]
    % Lecture note: Example 4.4.
    \label{lem:asym_example_interpolation}
    \lean{MultiBlockExamples.Running.exists_eval_zero_eq_zero_and_eval_eq,
        MultiBlockExamples.Running.eq_one_and_eq_one_of_forall_add_mul_pow_eq}
    \leanok
    In Example~4.4, let $\nu_0,\ldots,\nu_{m-1}$ be distinct nonzero complex numbers. For any
    $y_0,\ldots,y_{m-1}\in\C$, there is a polynomial $p$ with $p(0)=0$ and
    $p(\nu_j)=y_j$ for every $j$. In particular, if $c_0+c_1 2^n=1+2^n$ for
    every positive integer $n$, then $c_0=c_1=1$.
\end{lemma}

\begin{proof}
    \leanok
    Let $q$ be the Lagrange interpolant with $q(\nu_j)=y_j/\nu_j$, and put
    $p(X)=Xq(X)$. Since $\nu_j\neq0$, this gives $p(0)=0$ and
    $p(\nu_j)=\nu_j(y_j/\nu_j)=y_j$. For the last assertion the equations
    at $n=1,2$ already determine both coefficients.
\end{proof}

\subsection{A diagonal tensor with two different blocks}

\begin{lemma}[All reductions of a diagonal two-dimensional tensor]
    % Lecture note: Example 6.14.
    \label{lem:asym_example_diagonal_reductions}
    \lean{MultiBlockExamples.Running.tensorB,
        MultiBlockExamples.Running.blocks,
        MultiBlockExamples.Running.isInvariantSubmodule_tensorB_iff,
        MultiBlockExamples.Running.ord_flagVec_cases_of_z_eq_zero,
        MultiBlockExamples.Running.exists_multiBlockCompression_each_order,
        MultiBlockExamples.Running.exists_multiBlockCompression_upper,
        MultiBlockExamples.Running.exists_left_eq_upper_of_ord_lt,
        MultiBlockExamples.Running.exists_multiBlockCompression_swap,
        MultiBlockExamples.Running.exists_left_eq_swap_of_ord_gt,
        MultiBlockExamples.Running.redUpper_spec,
        MultiBlockExamples.Running.redSwap_spec}
    \leanok
    In Example~6.14, let $B=\operatorname{diag}(1,2)$ have one letter, and let the target blocks
    be $C_0=(1)$ and $C_1=(2)$. Its invariant subspaces are exactly
    $0,\C e_0,\C e_1,\C^2$. A compression onto these two blocks without zero
    slots has one of the two flags $0\subset\C e_0\subset\C^2$ and
    $0\subset\C e_1\subset\C^2$. In bond coordinates ordered by the flag, its
    gauge is respectively
    \begin{align}
        G            & =\begin{pmatrix}\alpha&\beta\\0&\delta\end{pmatrix}, &
        \widetilde G & =\begin{pmatrix}\beta&\alpha\\\delta&0\end{pmatrix},
                     & \alpha\delta\neq0.
    \end{align}
    Every such matrix defines a compression. For $G$, the reading rows are
    $W_0=(\alpha,\beta)$ and $W_1=(0,\delta)$, and the embedding columns are
    $V_0=(\alpha^{-1},0)^{\mathsf T}$ and
    $V_1=(-\beta/(\alpha\delta),\delta^{-1})^{\mathsf T}$.
    For $\widetilde G$, they are
    $\widetilde W_1=(\beta,\alpha)$, $\widetilde W_0=(\delta,0)$,
    $\widetilde V_1=(0,\alpha^{-1})^{\mathsf T}$ and
    $\widetilde V_0=(\delta^{-1},-\beta/(\alpha\delta))^{\mathsf T}$.
\end{lemma}

\begin{proof}
    \leanok
    If a subspace contains $(x,y)$, invariance under $B$ implies that it
    contains $(x,0)=2(x,y)-B(x,y)$ and $(0,y)=B(x,y)-(x,y)$.
    Thus its only possible one-dimensional members are the coordinate lines.
    The rows reading the first flag therefore have the first displayed form;
    reversing the order gives the second. Direct multiplication by the inverse
    gives upper triangular conjugated letters with the prescribed diagonal,
    and gives the displayed reading rows and embedding columns.
\end{proof}

\begin{lemma}[Reading maps that remain different after every word]
    % Lecture note: Examples 1.2 and 3.4.
    \label{lem:asym_example_diagonal_words}
    \lean{MultiBlockExamples.Running.exists_multiBlockCompression_gauge_one,
        MultiBlockExamples.Running.exists_multiBlockCompression_gauge_unipotent,
        MultiBlockExamples.Running.unipotent_conj_tensorB,
        MultiBlockExamples.Running.tensorB_pow,
        MultiBlockExamples.Running.left_mul_pow_ne_smul,
        MultiBlockExamples.Running.remainder_pow_ne_zero}
    \leanok
    In Examples~1.2 and~3.4, for the tensor of
    Lemma~\ref{lem:asym_example_diagonal_reductions}, the gauges $1$ and
    $H=\left(\begin{smallmatrix}1&1\\0&1\end{smallmatrix}\right)$
    have the same flag and satisfy
    $HBH^{-1}=\left(\begin{smallmatrix}1&1\\0&2\end{smallmatrix}\right)$.
    The reading rows of $C_0$ are $W=(1,0)$ and $\widetilde W=(1,1)$.
    For every $n\geq0$, $WB^n=(1,0)$ and $\widetilde WB^n=(1,2^n)$ are not
    proportional. Removing only $C_0$ leaves
    $B-e_0(1)e_0^{\mathsf T}=\operatorname{diag}(0,2)$, whose every
    nonnegative power is nonzero.
\end{lemma}

\begin{proof}
    \leanok
    Compute $B^n=\operatorname{diag}(1,2^n)$. The second coordinate of a scalar
    multiple of $(1,0)$ is zero, whereas $2^n\neq0$.
\end{proof}

\begin{lemma}[Matching and factorization in the diagonal example]
    % Lecture note: Example 6.14.
    \label{lem:asym_example_diagonal_matching}
    \lean{MultiBlockExamples.Running.matching_left_right_gOne_redSwap,
        MultiBlockExamples.Running.isAdaptedFactorization_gOne_redSwap}
    \leanok
    Compare $G=1$ with
    $\widetilde G=\left(\begin{smallmatrix}\beta&\alpha\\\delta&0\end{smallmatrix}\right)$,
    where $\alpha\delta\neq0$. The matching sends each labelled block to the
    same labelled block in the opposite position. For $C_0$ its intersection
    and lower sum are $U_0=\C e_0$ and $\mathcal N_0=\C e_1$; for $C_1$ they
    are $U_1=\C e_1$ and $\mathcal N_1=\C e_0$. The reading comparisons have
    scalars $\lambda_0=\delta$ and $\lambda_1=\alpha$, and
    \begin{align}
        \widetilde V_0-\delta^{-1}V_0
                       & =\begin{pmatrix}0\\-\beta/(\alpha\delta)\end{pmatrix}, &
        \widetilde V_1 & =\alpha^{-1}V_1.
    \end{align}
    With $L_0=\C e_0$, $L_1=\C e_1$, the adapted coordinates are $\Gamma=1$,
    and the factorization is
    \begin{align}
        \widetilde G G^{-1}
         & =\begin{pmatrix}1&\beta/\delta\\0&1\end{pmatrix}
        \begin{pmatrix}0&\alpha\\\delta&0\end{pmatrix}.
    \end{align}
\end{lemma}

\begin{proof}
    \leanok
    \uses{lem:asym_example_diagonal_reductions}
    The coordinate lines have zero intersection, so the matching reverses
    positions. Restrict the rows and columns from
    Lemma~\ref{lem:asym_example_diagonal_reductions} to these lines and multiply
    the two displayed factors.
\end{proof}

\subsection{Two copies of the same normal block}

\begin{lemma}[Multiplicity subspaces and split compressions]
    % Lecture note: Example 6.13.
    \label{lem:asym_example_multiplicity}
    \lean{MultiBlockExamples.Multiplicity.tensorC,
        MultiBlockExamples.Multiplicity.tensorB,
        MultiBlockExamples.Multiplicity.wordSpan_tensorC_three_eq_top_and_isNormal,
        MultiBlockExamples.Multiplicity.tensorB_eq_kronecker,
        MultiBlockExamples.Multiplicity.isInvariantSubmodule_lineTensor,
        MultiBlockExamples.Multiplicity.span_eq_of_lineTensor_eq,
        MultiBlockExamples.Multiplicity.commute_tensorB_iff_eq_kronMat,
        MultiBlockExamples.Multiplicity.exists_split_multiBlockCompression_kronRow,
        MultiBlockExamples.Multiplicity.exists_left_eq_kronRow_of_remainder_eq_zero}
    \leanok
    In Example~6.13, let $C=(\left(\begin{smallmatrix}0&1\\1&0\end{smallmatrix}\right),E_{00})$
    and $B^i=1_2\otimes C^i$. The words of length $3$ in $C$ span $\MN{2}$,
    so $C$ is normal. For every $u\in\C^2$, the subspace
    $\C u\otimes\C^2$ is invariant under $B$, and equality of two such
    subspaces for nonzero $u,\widetilde u$ implies
    $\C u=\C\widetilde u$. A matrix commutes with both letters of $B$ exactly
    when it is $A\otimes1_2$ for some $A\in\MN{2}$.

    The split compressions onto two copies of $C$ have no zero slots. They
    are exactly those with reading maps $W_j=a_j\otimes1_2$ and embedding
    maps $V_j=u_j\otimes1_2$, where $a_j$ are the rows of an invertible
    matrix $A$ and $u_j$ are the columns of $A^{-1}$. Every invertible $A$
    supplies such a split compression.
\end{lemma}

\begin{proof}
    \leanok
    The words $E_{00}^3$, $E_{00}^2\sigma_x$,
    $\sigma_xE_{00}^2$ and $\sigma_xE_{00}\sigma_x$ are the four matrix
    units. The identity $B^i(u\otimes v)=u\otimes C^iv$ proves invariance.
    Writing a commuting matrix as four $2\times2$ blocks makes each block a
    scalar. The same block equations identify all rectangular intertwiners
    as $a\otimes1_2$ and $u\otimes1_2$; the compression identities then give
    $AU=UA=1$.
\end{proof}

\begin{lemma}[Proportionality and matching with multiplicity]
    % Lecture note: Example 6.13.
    \label{lem:asym_example_multiplicity_matching}
    \lean{MultiBlockExamples.Multiplicity.exists_left_mul_evalWord_eq_smul_iff,
        MultiBlockExamples.Multiplicity.matching_left_splitRed_of_span_ne,
        MultiBlockExamples.Multiplicity.matching_left_redOne_redAt}
    \leanok
    For the tensor of Lemma~\ref{lem:asym_example_multiplicity}, suppose that
    two reading maps of the same labelled block are $a\otimes1_2$ and
    $\widetilde a\otimes1_2$. There exist a word $w$ with $C^w\neq0$ and a
    scalar $c$ such that
    $(\widetilde a\otimes1_2)B^w=c(a\otimes1_2)B^w$ if and only if
    $\widetilde a=ca$ for some scalar $c$. Each such scalar gives the identity
    for every word.

    Now let $A,\widetilde A$ define two split compressions, let $u_0$ be the
    first column of $A^{-1}$, and let $\widetilde u_0,\widetilde u_1$ be the
    columns of $\widetilde A^{-1}$. If $\C u_0\neq\C\widetilde u_0$, the
    matching pairs the first position of the first compression with the
    second position of the second. For this pair,
    $U=\C u_0\otimes\C^2$, $\mathcal N=\C\widetilde u_0\otimes\C^2$, and
    $\lambda=\widetilde a_1u_0\neq0$. On $U$ the second reading map is
    $\lambda$ times the first, and
    $\widetilde u_1-\lambda^{-1}u_0\in\C\widetilde u_0$.
    For $A=1$ and
    $\widetilde A=\left(\begin{smallmatrix}1&0\\1&1\end{smallmatrix}\right)$,
    these vectors are $u_0=e_0$, $\widetilde u_0=(1,-1)^{\mathsf T}$,
    $\widetilde u_1=e_1$, and $\lambda=1$.
\end{lemma}

\begin{proof}
    \leanok
    \uses{lem:asym_example_multiplicity}
    The word identity is
    $(\widetilde a-ca)\otimes C^w=0$; a nonzero entry of $C^w$ gives
    $\widetilde a=ca$. Distinct first lines give distinct first invariant
    subspaces with zero intersection. The inverse-matrix identities give
    $\widetilde a_1u_0\neq0$ and the stated difference of columns.
\end{proof}

\subsection{One block with two different flags}

\begin{lemma}[A split and a nonsplit compression onto one scalar block]
    % Lecture note: Example 6.15.
    \label{lem:asym_example_one_block}
    \lean{MultiBlockExamples.OneBlock.isInvariantSubmodule_span_tensorB_iff,
        MultiBlockExamples.OneBlock.z_eq_one_and_gInt_conj,
        MultiBlockExamples.OneBlock.flagVec_one_firstCompression_and_lastCompression,
        MultiBlockExamples.OneBlock.eq_toFlagData_firstCompression_or_lastCompression,
        MultiBlockExamples.OneBlock.left_right_matching_firstCompression_lastCompression,
        MultiBlockExamples.OneBlock.left_mulVec_e_one_and_left_mul_tensorB,
        MultiBlockExamples.OneBlock.remainder_lastCompression_eq_zero_and_firstCompression_ne_zero}
    \leanok
    In Example~6.15, let
    $B=\left(\begin{smallmatrix}1&1\\0&0\end{smallmatrix}\right)$ and
    $C=(1)$. The invariant lines of $B$ are $\C e_0$ and $\C(e_0-e_1)$;
    every labelled invariant flag for a compression onto $C$ is one of
    $0\subset\C e_0\subset\C^2$ and
    $0\subset\C(e_0-e_1)\subset\C^2$.
    The identity gauge puts $C$ first, whereas
    $\widetilde G=\left(\begin{smallmatrix}0&-1\\1&1\end{smallmatrix}\right)$
    satisfies $\widetilde GB\widetilde G^{-1}=\operatorname{diag}(0,1)$.
    Their matched reading maps and embeddings are
    $W=(1,0)$, $\widetilde W=(1,1)$, and $V=\widetilde V=e_0$.
    The matching sends the first block position to the second,
    $U=\C e_0$, and $\mathcal N=\C(e_0-e_1)$.
    Although $\widetilde We_1=1\neq0=We_1$, one has
    $\widetilde WB=WB$. The second compression is split and the first is not.
\end{lemma}

\begin{proof}
    \leanok
    Solving $Bv=\lambda v$ gives the two lines. The respective eigenvalues
    are $1$ and $0$, which determine their labels and hence the two possible
    flags. Multiplication by the displayed gauges and their inverses gives
    the rows, columns and remainders.
\end{proof}

\begin{lemma}[Sharpness for every number of zero slots]
    % Lecture note: Example 6.10.
    % The two sharpness conclusions are linked at cor:asym_uniq_one_block.
    \label{lem:asym_example_sharp_bound}
    \lean{MultiBlockExamples.SharpnessOneBlock.tensorBGen,
        MultiBlockExamples.SharpnessOneBlock.exists_multiBlockCompression_block_first,
        MultiBlockExamples.SharpnessOneBlock.exists_multiBlockCompression_block_last,
        MultiBlockExamples.SharpnessOneBlock.rowE0_mul_pow_and_rowOnes_mul_pow,
        MultiBlockExamples.SharpnessOneBlock.pow_mul_right_eq_of_eq_colE0}
    \leanok
    In Example~6.10, for $z\geq1$, let the one-letter tensor $B$ on $\C^{1+z}$ satisfy
    $Be_0=e_0$ and $Be_j=e_{j-1}$ for $1\leq j\leq z$, and put $C=(1)$.
    It admits compressions with $C$ first and last, with respective reading
    rows $W=e_0^{\mathsf T}$ and
    $\widetilde W=\sum_{j=0}^z e_j^{\mathsf T}$, and common embedding $e_0$.
    For every $n\geq0$,
    $WB^n=\sum_{j=0}^{\min(n,z)}e_j^{\mathsf T}$ and
    $\widetilde WB^n=\widetilde W$. These rows are proportional exactly when
    $n\geq z$. In particular, proportionality fails at the word of length
    $z-1=D_B-D_C-1$.
\end{lemma}

\begin{proof}
    \leanok
    The identity basis puts $C$ first. The basis
    $e_0-e_1,\ldots,e_{z-1}-e_z,e_0$ puts it last.
    Induction on $n$ gives the two row formulas; their zeroth coordinates
    force the proportionality scalar to be $1$.
\end{proof}

\subsection{A four-dimensional example with a two-dimensional block}

\begin{definition}[The two-letter four-dimensional tensor]
    % Lecture note: Example 6.16.
    \label{def:asym_example_worked}
    \lean{MultiBlockExamples.Worked.tensorB,
        MultiBlockExamples.Worked.tensorC₁,
        MultiBlockExamples.Worked.tensorC₂,
        MultiBlockExamples.Worked.tensorT,
        MultiBlockExamples.Worked.tensorC₁',
        MultiBlockExamples.Worked.matY,
        MultiBlockExamples.Worked.g0Int,
        MultiBlockExamples.Worked.gtInt}
    \leanok
    In Example~6.16, put
    \begin{align}
        B^0          & =\begin{pmatrix}-2&-2&-1&0\\2&3&1&0\\-2&-4&-1&0\\2&2&1&1\end{pmatrix}, &
        B^1          & =\begin{pmatrix}2&0&1&3\\0&2&0&-2\\-2&-2&-1&-1\\0&0&0&0\end{pmatrix},    \\
        G            & =\begin{pmatrix}2&2&1&1\\0&0&0&1\\1&1&1&1\\0&-1&0&1\end{pmatrix},      &
        \widetilde G & =\begin{pmatrix}3&2&2&3\\-2&-3&-1&1\\2&2&1&2\\0&0&0&1\end{pmatrix}.
    \end{align}
    Let $C_0=(\left(\begin{smallmatrix}0&1\\1&0\end{smallmatrix}\right),E_{00})$,
    $C_1=((1),(2))$,
    $T=(\left(\begin{smallmatrix}0&1\\0&1\end{smallmatrix}\right),
        \left(\begin{smallmatrix}0&0\\0&2\end{smallmatrix}\right))$, and
    $Y=\left(\begin{smallmatrix}1&1\\0&1\end{smallmatrix}\right)$.
    Write $C'_0=YC_0Y^{-1}$ and $C'_1=C_1$.
\end{definition}

\begin{lemma}[Reductions and traces of the four-dimensional tensor]
    % Lecture note: Example 6.16.
    \label{lem:asym_example_worked_reductions}
    \lean{MultiBlockExamples.Worked.wordSpan_tensorC₁_three_eq_top_and_isNormal,
        MultiBlockExamples.Worked.isInvariantSubmodule_span_tensorT_iff_and_mulVec_eq_zero,
        MultiBlockExamples.Worked.redG_spec,
        MultiBlockExamples.Worked.redGt_spec,
        MultiBlockExamples.Worked.tensorC₁'_eq,
        MultiBlockExamples.Worked.isNormal_blocks,
        MultiBlockExamples.Worked.isNormal_blocks',
        MultiBlockExamples.Worked.trace_evalWord_tensorB_eq_add}
    \leanok
    \uses{def:asym_example_worked}
    Both families $(C_0,C_1)$ and $(C'_0,C'_1)$ are normal. The tensor $T$
    has the unique invariant line $\C e_0$, and both letters vanish on it.
    The gauge $G$ gives $GB^iG^{-1}=C_0^i\oplus T^i$ and is a compression
    in the order $(C_0,0,C_1)$. The gauge $\widetilde G$ gives a compression
    in the order $(0,C'_1,C'_0)$, with conjugated letters
    \begin{align}
        \widetilde GB^0\widetilde G^{-1}
         & =\begin{pmatrix}0&2&2&-3\\0&1&2&-4\\0&0&1&0\\0&0&1&-1\end{pmatrix}, &
        \widetilde GB^1\widetilde G^{-1}
         & =\begin{pmatrix}0&2&3&-5\\0&2&1&-3\\0&0&1&-1\\0&0&0&0\end{pmatrix}.
    \end{align}
    For every nonempty word $w$,
    $\tr(B^w)=\tr(C_0^w)+\tr(C_1^w)$; the trace at the empty word is $4$.
\end{lemma}

\begin{proof}
    \leanok
    Normality of $C_0$ follows from its length-three words spanning the matrix
    units, and $C'_0$ is its gauge transform. The scalar block is nonzero.
    The two equations $T^iv=\lambda_i v$ force $v\in\C e_0$.
    Direct multiplication gives the conjugated letters. In a nonempty word,
    the zero diagonal slot contributes zero to the trace.
\end{proof}

\begin{lemma}[Flags and comparison maps of the four-dimensional tensor]
    % Lecture note: Example 6.16.
    \label{lem:asym_example_worked_flags}
    \lean{MultiBlockExamples.Worked.flagVec_redG_redGt_and_matching,
        MultiBlockExamples.Worked.inf_sup_intertwiner_redG_redGt,
        MultiBlockExamples.Worked.not_exists_left_redGt_eq_smul_left_redG,
        MultiBlockExamples.Worked.finrank_flagVec_redG_inf_flagVec_redGt}
    \leanok
    \uses{lem:asym_example_worked_reductions}
    Put $a=(1,0,-1,0)^{\mathsf T}$, $b=(-1,1,-1,1)^{\mathsf T}$,
    $c=(-1,0,2,0)^{\mathsf T}$ and $d=(2,-1,-2,0)^{\mathsf T}$.
    The two flags have
    $M_1=\spn_{\C}\{a,b\}$, $M_2=\spn_{\C}\{a,b,c\}$,
    $\widetilde M_1=\C c$, $\widetilde M_2=\spn_{\C}\{c,d\}$, and
    $M_3=\widetilde M_3=\C^4$. Their intersection dimensions are
    \begin{align}
        \bigl(\dim(M_k\cap\widetilde M_{k'})\bigr)_{1\leq k,k'\leq3}
         & =\begin{pmatrix}0&0&2\\1&1&3\\1&2&4\end{pmatrix}.
    \end{align}
    The matching sends the positions $(1,2,3)$ to $(3,1,2)$.
    For the $C_0$ pair, $U=\spn_{\C}\{a,(-2,1,0,1)^{\mathsf T}\}$,
    $K=0$, and $\mathcal N=\spn_{\C}\{c,d\}$; its reading comparison is
    $\widetilde W_0x=YW_0x$ on $U$, and every column of
    $\widetilde V_0-V_0Y^{-1}$ lies in $\mathcal N$.
    For the zero pair the readings agree on $U$ and $\dim K=0$; for the
    $C_1$ pair they agree on $U$ and $\dim K=1$.
    The full reading rows of $C_1$ are
    $W_1=(0,-1,0,1)$ and $\widetilde W_1=(-2,-3,-1,1)$, which are not proportional.
\end{lemma}

\begin{proof}
    \leanok
    The vanishing coordinates of $Gx$ and $\widetilde Gx$ give the four
    subspaces. Their intersections determine the matching and the dimension
    matrix. Restrict the rows of the two gauges to the stated intersections;
    the columns of the inverse gauges give the comparison modulo the lower
    sums. The first coordinate rules out proportionality of $W_1$ and
    $\widetilde W_1$.
\end{proof}

\begin{lemma}[An explicit basis adapted to the two four-dimensional flags]
    % Lecture note: Example 6.16.
    \label{lem:asym_example_worked_factorization}
    \lean{MultiBlockExamples.Worked.complementVec,
        MultiBlockExamples.Worked.complement,
        MultiBlockExamples.Worked.gammaGauge,
        MultiBlockExamples.Worked.isAdaptedFactorization_redG_redGt}
    \leanok
    \uses{lem:asym_example_worked_flags}
    In Example~6.16, with the notation of Lemma~\ref{lem:asym_example_worked_flags}, choose
    $L_1=M_1$, $L_2=\C c$ and $L_3=\C(0,-1,2,0)^{\mathsf T}$.
    The adapted coordinates and factors are
    \begin{align}
        \Gamma                & =\begin{pmatrix}
                                     2 & 2  & 1 & 1 \\
                                     0 & 0  & 0 & 1 \\
                                     1 & 2  & 1 & 0 \\
                                     0 & -1 & 0 & 1
                                 \end{pmatrix},  &
        \mathsf A             & =1+E_{23},         \\
        \widetilde{\mathsf A} & =\begin{pmatrix}
                                     1 & 2 & 1  & -1 \\
                                     0 & 1 & -1 & 2  \\
                                     0 & 0 & 1  & 0  \\
                                     0 & 0 & 0  & 1
                                 \end{pmatrix}, &
        P_X                   & =\begin{pmatrix}
                                     0 & 0 & 1 & 0 \\
                                     0 & 0 & 0 & 1 \\
                                     1 & 1 & 0 & 0 \\
                                     0 & 1 & 0 & 0
                                 \end{pmatrix}.
    \end{align}
    Here $X_{C_0}=Y$ and $X_0=X_{C_1}=1$. The factors $\mathsf A$ and
    $\widetilde{\mathsf A}$ are block upper triangular with identity diagonal
    blocks in the two flag orders, and
    $G=\mathsf A\Gamma$,
    $\widetilde G=\widetilde{\mathsf A}P_X\Gamma$, and
    $\widetilde GG^{-1}=\widetilde{\mathsf A}P_X\mathsf A^{-1}$.
\end{lemma}

\begin{proof}
    \leanok
    The columns of $\Gamma^{-1}$ are bases of the three chosen complements.
    Matrix multiplication gives the two gauge identities, and the displayed
    diagonal blocks are identities.
\end{proof}

\subsection{Reading flags and explicit factorizations in coordinates}

\begin{lemma}[Flags, reading rows and embedding columns of an explicit gauge]
    \label{lem:asym_example_flag_coordinates}
    \lean{MultiBlockExamples.flagVec,
        MultiBlockExamples.mem_flagVec_iff,
        MultiBlockExamples.mem_flagVec_gaugeOfMatrix,
        MultiBlockExamples.leftRow_gaugeOfMatrix,
        MultiBlockExamples.rightCol_gaugeOfMatrix,
        MultiBlockExamples.matching_eq_of_flagVec_inf}
    \leanok
    Let a compression have gauge matrix $G$, and index the rows of $G$ by
    its block coordinates. Its flag is
    $M_k=\{x:(Gx)_b=0\text{ whenever }\sigma(b)\geq k\}$, where positions
    are counted from $0$. Its reading map $W_b$ is the block row of $G$,
    and its embedding map $V_b$ is the corresponding block column of $G^{-1}$.
    For two compressions onto normal blocks of positive bond dimension,
    the matching of the step at position $k$ is the least $j$ such that
    $M_{k+1}\cap\widetilde M_{j+1}$ is not contained in $M_k$.
\end{lemma}

\begin{proof}
    \leanok
    Apply the gauge to the leading coordinate subspaces. The reading and
    embedding maps are the coordinate projection and inclusion transported
    through $G$ and $G^{-1}$. The characterization of the matching is its
    first nonzero step on the quotient $M_{k+1}/M_k$.
\end{proof}

\begin{definition}[An explicit adapted factorization]
    \label{def:asym_example_adapted_factorization}
    \lean{MultiBlockExamples.IsAdaptedFactorization}
    \leanok
    For two compressions onto normal blocks of positive bond dimension,
    an adapted factorization consists of their matching on labelled slots,
    complements $L_k$ of $K_k$ in $U_k$, an invertible coordinate map
    $\Gamma$ whose restriction to $L_k$ reads the corresponding block through
    $W_{\sigma(k)}$, invertible matched-block intertwiners $X_b$, and two
    block upper triangular matrices $\mathsf A,\widetilde{\mathsf A}$ with
    identity diagonal blocks. They satisfy
    $\widetilde W_{\pi(b)}x=X_bW_bx$ on the matched intersection and
    $G=\mathsf A\Gamma$ and
    $\widetilde G=\widetilde{\mathsf A}P_X\Gamma$, where
    $P_X=\sum_b\iota_{\pi(b)}X_bp_b$, and $p_b,\iota_b$ are the coordinate
    block projection and inclusion.
\end{definition}
